\section{Related Work}\label{sec:related}

\paragraph{Activation functions.}
ReLU~\citep{nair2010rectified} enabled deep training; its piecewise-linear form provides scale invariance but only binary gating.
ELU~\citep{clevert2016fast} and SELU~\citep{klambauer2017self} introduced smooth negative tails but are not self-gated.
SiLU/Swish~\citep{elfwing2018sigmoid,ramachandran2017searching} and GELU~\citep{hendrycks2016gaussian} popularized the self-gated $z \cdot g(z)$ form; Mish~\citep{misra2019mish} is a further variant.
All share the scale sensitivity analyzed in Section~\ref{sec:tradeoff}.
NELU is, to our knowledge, the first self-gated activation to restore it.

\paragraph{Gated linear units.}
SwiGLU~\citep{shazeer2020glu} and GeGLU use a separate projection as gate, partially decoupling scale from gating at the cost of doubling gate parameters.
NELU achieves the decoupling within the self-gated mechanism itself, without additional parameters.

\paragraph{Normalization.}
BatchNorm~\citep{ioffe2015batch}, LayerNorm~\citep{ba2016layer}, RMSNorm~\citep{zhang2019root}, and GroupNorm~\citep{wu2018group} normalize \emph{between} layers.
NELU normalizes \emph{inside} the gate, targeting scale sensitivity that inter-layer normalization cannot reach (Section~\ref{sec:consequences}).

\paragraph{Combined normalization-activation.}
EvoNorm-S0~\citep{liu2020evolving}, discovered via evolutionary search, has the form $x \cdot \sigma(v_1 x)\,/\,\mathrm{std}_{\mathrm{group}}(x)$.
Despite structural similarity, it is \emph{not} scale-invariant: the gate argument $v_1 x$ still depends on absolute scale, and the output division removes magnitude information entirely.
NELU normalizes only the gate input, yielding provable scale invariance ($\mathrm{NELU}(\alpha\mathbf{z}) = \alpha\,\mathrm{NELU}(\mathbf{z})$) while preserving output magnitude---a property EvoNorm-S0 lacks.
The EvoNorm authors observe post-hoc that their discovered layers ``attempt to promote scale-invariance''; our work provides the theoretical analysis of \emph{why} it matters and the minimal fix that provably achieves it.

\paragraph{Scale invariance in weights.}
Weight normalization~\citep{salimans2016weight} and spectral normalization~\citep{miyato2018spectral} constrain the weight matrix.
\citet{li2022robust} showed that making the \emph{architecture} scale-invariant enables robust training with SGD alone; NELU achieves scale invariance in the \emph{activation}, a complementary and composable intervention.


\section{Conclusion}\label{sec:conclusion}

Smooth self-gated activations gained graded modulation and differentiable gradients over ReLU, but lost scale invariance: their output direction---the signal downstream layers use for feature selection---becomes entangled with pre-activation scale.
We proved this is inherent to any non-constant smooth gate in the $z \cdot g(z)$ form, and showed that weight growth during training drives the gate toward saturation, progressively reducing the graded-gating regime.

NELU restores scale invariance by normalizing the gate's input via stop-gradient RMS.
Scale is not discarded but redirected: the gate sees only relative magnitudes; absolute scale flows through the output magnitude.
NELU is, to our knowledge, the first smooth self-gated activation with provably scale-invariant output direction, combining GELU's differentiability with ReLU's immunity to scale.
Across architectures and domains---CNNs on CIFAR, ViTs on ImageNet, and GPT-2 language models up to 774M parameters---NELU matches or improves GELU on clean benchmarks and substantially improves robustness to corruption, with the largest gains where pre-activation scale is most disrupted.

\paragraph{Limitations.}
We evaluate only the GELU variant ($g = \Phi$); the SiLU variant ($z_i \cdot \sigma(z_i / \operatorname{sg}[\rho])$, Section~\ref{sec:definition}) enjoys identical scale-invariance guarantees but is not benchmarked here.
The optimal normalization axis differs between CNNs and Transformers (Appendix~\ref{app:ablation}); a unified rule remains future work.
The link from scale sensitivity to robustness is supported by consistent experimental evidence, not formal proof.